\documentclass[a4paper, 12pt]{report}
\usepackage[utf8]{inputenc}    
\usepackage[T1]{fontenc}
\usepackage[french]{babel}
\usepackage{hyperref}
\usepackage{setspace}
\usepackage{newtxtext,newtxmath}
\usepackage[top=3cm, bottom=3cm, left=3cm, right=3cm]{geometry}
\usepackage[explicit]{titlesec}
\usepackage{graphicx}
\usepackage[stable]{footmisc}
\usepackage{wrapfig}
\usepackage{multicol}
\usepackage{minitoc}
\usepackage{plantuml}

%* Config

\titleformat{\chapter}[display]{}{}{0pt} {
  \parbox{\textwidth} {
    \textbf{\fontsize{30}{30}\selectfont\thechapter.\hspace{0.5cm}\LARGE#1}\\
    \rule{\textwidth}{0.4pt}
  }
}
\titleformat{name=\chapter,numberless}[display]{}{}{0pt} {
  \parbox{\textwidth}{
    \LARGE\textbf{#1}\\
    \rule{\textwidth}{0.4pt}
  }
}

\hypersetup{
    colorlinks=true,
    urlcolor=black,
    linkcolor= black,
    citecolor=black
}

\mtcsetrules{*}{off}

\begin{document}
\renewcommand \partname{\thispagestyle{empty}}
\doparttoc

%* Page de garde

\thispagestyle{empty}
\begin{center}  
\begin{multicols}{2}
  \flushleft
  \null
  \includegraphics[width=\columnwidth]{../res/logo-iut.png}
  \flushright
  \null
  \vspace{0.3cm}
  \large{Promotion 2020-2021}
\end{multicols}
\vspace{2cm}
\LARGE{\textit{Mise en place de traitements massifs et asynchrones dans le cadre de l'import de relèves}}\\
\vspace{2cm}
\large{\textbf{Mémoire présenté en septembre 2021 par}}\\
\vspace{0.5cm}
\Large{\textit{Loïc STEINMETZ}}\\
\vspace{2cm}
\large{\textbf{En vue de l'obtention de la Licence Professionnelle}}\\
\vspace{0.5cm}
\fbox{
  \begin{minipage}{12cm}
    \begin{center}
      \null
      \vspace{0.3cm}
      \textbf{Métiers de l'Informatique}\\
      Conception, développement et test de logiciels\\
      Parcours Métiers du Génie Logiciel\\
      \null
    \end{center}
  \end{minipage}
}
\vspace{2.8cm}
\begin{multicols}{2}
  \flushleft
  \null
  \textbf{Alternance effectuée à :}\\
  Efluid\\
  2 bis, rue Ardant du Picq\\
  CS 10100 57004 METZ Cedex 01\\
  \flushright
  \null
  \includegraphics[width=.6\columnwidth]{../res/logo-efluid.jpg}
\end{multicols}
\end{center}

\clearpage
\thispagestyle{empty}
\null
\clearpage

\thispagestyle{empty}
\begin{center}  
\begin{multicols}{2}
  \flushleft
  \null
  \includegraphics[width=\columnwidth]{../res/logo-iut.png}
  \flushright
  \null
  \vspace{0.3cm}
  \large{Promotion 2020-2021}
\end{multicols}
\vspace{2cm}
\LARGE{\textit{Mise en place de traitements massifs et asynchrones dans le cadre de l'import de relèves}}\\
\vspace{2cm}
\large{\textbf{Mémoire présenté en septembre 2021 par}}\\
\vspace{0.5cm}
\Large{\textit{Loïc STEINMETZ}}\\
\vspace{2cm}
\large{\textbf{En vue de l'obtention de la Licence Professionnelle}}\\
\vspace{0.5cm}
\fbox{
  \begin{minipage}{12cm}
    \begin{center}
      \null
      \vspace{0.3cm}
      \textbf{Métiers de l'Informatique}\\
      Conception, développement et test de logiciels\\
      Parcours Métiers du Génie Logiciel\\
      \null
    \end{center}
  \end{minipage}
}
\vspace{2.8cm}
\begin{multicols}{2}
  \flushleft
  \null
  \textbf{Alternance effectuée à :}\\
  Efluid\\
  2 bis, rue Ardant du Picq\\
  CS 10100 57004 METZ Cedex 01\\
  \flushright
  \null
  \includegraphics[width=.6\columnwidth]{../res/logo-efluid.jpg}
\end{multicols}
\end{center}

%* Remerciements

\chapter*{Remerciements}
\addcontentsline{toc}{chapter}{Remerciements}
\thispagestyle{empty}

Je remercie tout particulièrement mon maître de stage, Alexandre L'HUILLIER, pour m'avoir fait bénéficier de ses compétences tout au long de mon année d'alternance. Merci à lui pour son professionalisme, sa confiance et sa disponibilité.\\

Je remercie également Jean-Luc THIRY, chef de service et Didier GRZEJSZCZAK, chef de filière, pour leur accompagnement général et pour ma bonne intégration au sein de l'entreprise.\\

Mes remerciements vont enfin à l'ensemble des membres du service auquel j'ai été intégré et que je n'ai pas encore cités : Patrice FRANTZ, Julien KEMPF, Fabienne MAMET, François MOLIN, Roderick PIERRE, Stéphane POIROT et Stéphane ROEDEL.

%* Table des matières : mémoire

\part{Mémoire}
\renewcommand{\clearpage}{}
\chapter*{Sommaire}
\renewcommand\ptctitle{}
\parttoc
\thispagestyle{empty}
\renewcommand{\clearpage}{\newpage}
\clearpage

%* Abstract

\chapter*{Abstract}
\addcontentsline{toc}{chapter}{Abstract}

%* Introduction

\chapter*{Introduction}
\addcontentsline{toc}{chapter}{Introduction}

%* Présentation de l'entreprise

\chapter{Présentation de l'entreprise}

%* Sujets traités

\chapter{Sujets traités}

\section{Généralités}
\section{Développements courants}
\section{Projet principal}

%* Traitement des sujets

\chapter{Traitement du projet industriel}

\section{Analyses préliminaires}
\section{Formalisation du besoin et planification}
\section{Conception techniques}
\section{Développements fonctionnels}
\section{Tests}

%* Conclusion

\chapter*{Conclusion}
\addcontentsline{toc}{chapter}{Conclusion}

%* Table des matières : annexes

\part{Annexes}
\renewcommand{\clearpage}{}
\chapter*{Table des annexes}
\renewcommand\ptctitle{}
\parttoc
\thispagestyle{empty}
\renewcommand{\clearpage}{\newpage}
\appendix

%* Annexes

\chapter{Efluid : répartition des parts}
\label{appendix:efluid-parts}

\begin{center}
  \begin{plantuml}
    @startuml
    skinparam monochrome true
    rectangle UEM
    rectangle CDC
    rectangle Enedis
    rectangle Efluid
    UEM -- Efluid : 60%
    CDC -- Efluid : 10%
    Enedis -- Efluid : 30%
    @enduml
  \end{plantuml}
\end{center}

\vspace{1cm}
\noindent\textbf{UEM :} Usine d'Electricité de Metz\\
\noindent\textbf{CDC :} Caisse des Dépôts et Consignations

\chapter{Efluid : organisation des services}
\label{appendix:efluid-organisation}

\begin{center}
  \begin{plantuml}
    @startuml
    skinparam monochrome true
    left to right direction
    rectangle Présidence as P
    rectangle "Direction informatique" as DI
    rectangle "Direction commerciale" as DC
    rectangle Ginko as G
    rectangle Technologies as T
    rectangle "Développement : CRM et Facturation" as DFac
    rectangle "Développement : Comptage et réseaux" as DCpt #BBB
    rectangle "Expertise fonctionnelle : CRM et Facturation" as EFac
    rectangle "Expertise fonctionnelle : Comptage et réseaux" as ECpt
    rectangle "Coordination et clients" as CC
    rectangle "Prestation et projets" as PP
    P -- DC
    P -- DI
    P -- G
    DI -- T
    DI -- DFac
    DI -- DCpt
    DC -- EFac
    DC -- ECpt
    DC -- CC
    DC -- PP
    @enduml
  \end{plantuml}
\end{center}

%* Tables des matières : général

\tableofcontents
\thispagestyle{empty}

%* Résumé

\chapter*{Résumé}
\thispagestyle{empty}

Résumé.

\end{document}